\section{Die Kapazität}

Als Kapazität wird das Verhältnis der Ladung $q$, die sich auf einem Leiter befindet, zu seiner elektrischen Spannung $U$ gegenüber einem Bezugspunkt definiert:
\[
    \text{Spannung:} \ U = \phi_B - \phi_A
\]\[
    \text{Kapazität:} \ C = \frac{q}{U}
\]
Einheit [C]: = Farad = F = $\frac{C}{V}$ = $\frac{Coulomb}{Volt}$

\subsection{Kondensatoren}
Ist eine Anordung von zwei Leiter, die beide eine gleich grosse, aber entgegengesetzte Ladung tragen können.\\[1ex]
Ladungen sind: $+q$ und $-q$ \\[2ex]
\textbf{Kapazität eines Kondensators:} 
\[
    C = \frac{q}{U} \ ; \ q \ \text{als Ladung auf einem Leiter und} \ U \ \text{als Spannung zwischen den Leiter}
\]

\textbf{Der Plattenkondensator}
Es gilt:
\[
    U = E \cdot d = \frac{\sigma}{\epsilon_0} \cdot d = \frac{q \cdot d}{\epsilon_0 \cdot A}
\]
$E = \frac{\sigma}{\epsilon_0}$: als $\Vec{E}-Feld$ zwischen den Oberflächen \\[1ex]
$d$: als Abstand zwischen den Platten \\[1ex]
$\sigma = \frac{q}{A}$: als Flächenladungdichte, $A$: Fläche der Platte \\[2ex]
\textbf{Kapazität eines Plattenkondensators:}
\begin{formulabox}
\[
     C = \frac{q}{U} = \frac{q}{(\frac{q \cdot d}{\epsilon_0 A})} = \frac{\epsilon_0 A}{d}
\]
\end{formulabox}
\textbf{Kapazität eines Zylinderkondensators:}
\begin{formulabox}
    \[
    C = \frac{q}{U} = \frac{2 \pi \epsilon_0 l}{\ln(\frac{r_2}{r_1})}
    \]
\end{formulabox}

\subsection{Speichern von elektrischer Energie}
Arbeit wird während dem Ladevorgang eines Kondensators verrichtet und kann als elektrische Energie gespeichert werden. \\[1ex]
Der Gesamtzuwachs an Energie wird wiefolgt beschrieben: 
\[
    E_{el} = \int_0^q dE_{el} = \int_0^q \frac{q'}{C}dq' = \frac{1}{C}\int_0^q q' dq' = \frac{1}{2}\frac{q^2}{C}
\]
\textbf{Die im Kondensator gespeicherte Energie:}
\begin{formulabox}
    \[
    \Rightarrow E_{el} = \frac{1}{2}\frac{q^2}{C} = \frac{1}{2}qU = \frac{1}{2}CU^2
    \]
\end{formulabox}

Definiert man $w_{el}$ als Energiedichte (Energie pro Volumen) so folgt: \\[1ex]
\textbf{Energiedichte des elektrischen Felds:}
\begin{formulabox}
    \[
        w_{el} = \frac{\text{Energie}}{\text{Volumen}} = \frac{1}{2}\epsilon_0 E^2
    \]
\end{formulabox}

\subsection{Kondensatoren, Batterien \& elektrische Stromkreise}
Über alle Parallelgeschaltenen Bauelementen in elektr. Stromkreisen, herrt die gleiche Spannung. \\[1ex]
Für die gespeicherte Gesamtladung gilt:
\[
    q = q_1 + q_2 = C_1 U + C_2 U = U(C_1 + C_2)
\]
Ein einzelner Kondensator mit gleicher Wirkung wie die zusammen geschaltenen nennt man Ersatzkapazität:\\[3ex]
Für \textbf{Parallelschaltung} gilt:
\[
    C = \frac{q}{U} = \frac{q_1 + q_2}{U} = \frac{q_1}{U} + \frac{q_2}{U} = C_1 + C_2
\]
\textbf{Ersatzkapazität für parallel geschaltete Kondensatoren}
\begin{formulabox}
    \[
        \Rightarrow C_{eq} = C_1 + C_2 + ... + C_i
    \]
\end{formulabox}

Für \textbf{Reihen-/ Serienschaltung} gilt:
\[
    U = \phi_A - \phi_B = U_1 + U_2 = \frac{q}{C_1} + \frac{q}{C_2} = q\left( \frac{1}{C_1} + \frac{1}{C_2} \right)
\]
\textbf{Ersatzkapazität für in Reihe geschaltete Kondensatoren}
\begin{formulabox}
    \[
        \Rightarrow C_{eq} = \left( \frac{1}{C_1} + \frac{1}{C_2} + ... + \frac{1}{C_i} \right)^{-1}
    \]
\end{formulabox}

\subsection{Dielektrika}
Dielektrikas sind nichtleitende Matrielien (z.B Holz, Glas, Plexiglas, etc.) \\[2ex]

\textbf{Kapazitäterhöhung} durch ein Dielektrikum: \\[1ex]
$\Rightarrow \Vec{E} = \frac{\Vec{E}_0}{\varepsilon_{rel}}$ \\[1ex]
Dielektrikum schwächt $\Vec{E}-Feld \rightarrow$ gleiche Ladung, kleinere Spannung $\rightarrow$ höhere Kapazität. \\[1ex]
$\varepsilon_{rel}$: hängt vom Dielektrikum ab, ist eine Matrielkonstante \\[2ex]
\begin{formulabox}
    \[
        \Rightarrow U = E \cdot d = \frac{E_0 \cdot d}{\varepsilon_{rel}} = \frac{U_0}{\varepsilon_{rel}}
    \]  
\end{formulabox}
\begin{formulabox}
    \[
        \Rightarrow C = \frac{q}{U} = \varepsilon_{rel}\frac{q}{U_0} = \varepsilon_{rel} \cdot C_0
    \]
\end{formulabox}

\textbf{gespeicherte Energie und die Energiedichte:}
\begin{formulabox}
    \[
        E_{el} = \frac{1}{2}CU^2 = \frac{1}{2}\varepsilon_{rel}\epsilon_0 E^2 A d
    \]
    \[
    w_{el} = \frac{1}{2}\varepsilon_{rel}\epsilon_0 E^2
    \]
\end{formulabox}